\documentclass[10pt,a4paper]{amsart}
\usepackage[margin=19mm]{geometry}
\usepackage{amsmath,amssymb,mathtools}
\usepackage[scr=rsfs,cal=euler]{mathalfa}
\usepackage{xfrac}
\usepackage[unicode,hidelinks,bookmarks=false]{hyperref}
\newcommand{\ie}{i.e.\ }
\newcommand{\cf}{cf.\ }
\newcommand{\resp}{respectively\ }
\newcommand{\Subsection}{\subsection}
\providecommand{\nobreakdash}{\nobreak\hbox{-}}
\setlength{\emergencystretch}{2em}
\multlinegap0pt
\usepackage{amssymb}
\usepackage{xcolor}
\newtheorem{theorem}{Theorem}
\def\r{\mathbb R}
\title{Newton's problem of minimal resistance in Lorentz-Minkowski space}
\author{Rafael L\'opez}
\date{Source: arXiv:2605.17096v1 (CC BY 4.0)}
\begin{document}
\maketitle

\noindent\emph{Source and licence.} Newton's problem of minimal resistance in Lorentz-Minkowski space. Version arXiv:2605.17096v1 (CC BY 4.0), distributed by arXiv under the Creative Commons Attribution 4.0 International licence, http://creativecommons.org/licenses/by/4.0/, as recorded in the arXiv licence metadata for this exact version and shown on the abstract page https://arxiv.org/abs/2605.17096v1. Full licence notice: \texttt{licenses/arxiv-2605.17096-CC-BY-4.0.txt}.\par\smallskip

\noindent\textit{Editorial scope.} Counted unit: the article's theorem t51 (source0.tex lines 427--441). The article's complete original proof of it is delivered with it (source0.tex lines 443--498, one \texttt{proof} environment of 56 lines). No mathematics is written, repaired or re-derived: the statement and the proof are the article's own lines, in the article's own notation, and no retained line is substituted or re-flowed.  Retained uncounted as the source-local context the proof needs: lines 198--200; lines 397--423.  Nothing was omitted from any retained span, so the package carries no omission record.  No retained line contains a citation, so the wrapper carries no bibliography and no bibliography entry.  No retained line contains a figure environment, an image-inclusion command or drawing code, so no image is delivered and the package declares no asset dependency.  Editorial formatting only, in addition: an AMS \texttt{amsart} preamble that restates the article's own declarations and macros used by the retained lines, namely the environments \texttt{theorem} on separate counters, together with the standard packages those lines need.  The retained environments print in this document as Theorem 1 in order of appearance, whereas the article prints the same items with the numbering of its own full document.  The complete native source of the article as distributed in the arXiv e-print for this exact version is bundled unchanged as \texttt{sources/200765/source0.tex}.
\begin{equation}\label{eqs}
(1-|Du|^2)\Delta u+4 D^2u(Du,Du)=0,
\end{equation}

We study the radial solutions of the Euler-Lagrange equation \eqref{eqs}, or equivalently,   stationary surfaces  which are surfaces of revolution about the $z$-axis. Since $S$ is spacelike, then $S$ is a locally graph over $\r^2$. If, for instance, we are looking for rotational spacelike surfaces defined on $B_R$, then $u$ is a global graph on $B_R$, and thus, $u$ can be written as $u=u(r)$. The surface $S$   can be described by its generating curve $\gamma(r)=(r,0,u(r))$, where 
$r\in [a,b]$, with  $a\geq 0$. Hence, $S=\{(r\cos t, r\sin t,u(r))\colon a\leq r\leq b,t\in\r\}$. The Lagrangian of the energy $E$ is 
\begin{equation}\label{ff}
F(r,u,u')=\frac{r}{1-u'^2}.
\end{equation}
Since $F$ does not depend explicitly on $u$, the Euler-Lagrange equation is  
$$\frac{d}{dr}\frac{ru'(r)}{(1-u'(r)^2)^2}=0.$$
Therefore, there exists  a constant $c_1\in\r$ such that 
\begin{equation}\label{r1}
\frac{ru'(r)}{(1-u'(r)^2)^2}=c_1.
\end{equation}
 If $c_1=0$, then $u(r)=c$ for some constant $c\in\r$, and $S$ is a horizontal spacelike plane. Suppose now $c_1\not=0$. Without loss of generality, from now on we assume $c_1>0$.  Equivalently, we have $u'>0$ in its entire domain due to \eqref{r1}. This proves that $u(r)$ is strictly increasing in its domain. Let $p=u'$ be the new independent variable.   Then \eqref{r1} becomes
\begin{equation}\label{r2}
r= \frac{c_1}{p}(1-p^2)^2.
\end{equation}
In particular, $r(p)$ is a decreasing function of $p$. 
We now find a parametric representation of $\gamma$ in terms of the variable $p$. Differentiating   \eqref{r2} gives 
$$\frac{dr}{dp}=c_1 (-\frac{1}{p^2}-2+3p^2).$$
Then
$$u=\int p\, dr= \int c_1 p\left(-\frac{1}{p^2}-2+3p^2\right)dp=c_2- c_1 \left(\log(p)+p^2-\frac{3p^4}{4}\right),$$
where we use  that fact that $p>0$. 
Thus, we  have $\gamma(p)=(r(p),u(p))$, where 
\begin{eqnarray}
r(p)&=&\frac{c_1}{p}(1-p^2)^2,\label{p1}\\
u(p)&=&c_2-c_1\left(\log(p)+p^2-\frac{3p^4}{4}\right).\label{p2}
\end{eqnarray}
We study the existence of radial extremals of $E$ parametrized by \eqref{p1}-\eqref{p2}. The first examples are horizontal planar disks $u(r)=c$, $c\in\r$. 

\begin{theorem}\label{t51}
 Let $M$ and $R$ be positive numbers with $M<R$. Then there is a unique extremal of $E$ of radial type, $u=u(r)$, given as a solution $u \colon (0,R]\to\r$  of \eqref{r1} such that 
\begin{equation}\label{c1}\lim_{r\to 0}u(r)=0,\quad u(R)=M.
\end{equation}
Moreover: 
\begin{enumerate}
\item  The function $u$ is strictly increasing and $0< u(r)\leq M$ for all $r\in (0,R]$.
\item The function $u$ can be extended to   $r=0$, with the property 
$$\lim_{r\to 0}u'(r)=1.$$
\item The curve $\gamma $ is regular   on its  domain.
\item The function $u(r)$ can be extended to $(0,\infty)$ while maintaining  the spacelike condition $u'^2<1$. Moreover, 
$$\lim_{r\to \infty}u(r)=\infty,\quad  \lim_{r\to \infty}u'(r)=0.$$
\item The function $u(r)$ is concave. 
\end{enumerate} 
\end{theorem}

\begin{proof}
Based on \eqref{r1} and \eqref{r2}, the function $u$ can be defined with the variable $p=u'$ according to \eqref{p1} and \eqref{p2}. We find the values of $c_1$ and $c_2$ that satisfy the boundary conditions \eqref{c1}. 

From \eqref{p1}, $u(r)$ is defined up to  $r=0$ as $p\to 1$. If we require $u(0)=0$,  from \eqref{p2} we deduce $c_2=c_1/4$. Replacing this in \eqref{p2}, we obtain 
\begin{equation}\label{p3}
u(p)=c_1\left(\frac14-\log(p)-p^2+\frac{3p^4}{4}\right).
\end{equation}
We now find the value of $p$ such that $u(R)=M$, which  is the second condition in \eqref{c1}. The function   $r\colon (0,1]\to\r$, $r=r(p)$, is  strictly decreasing with 
$$\lim_{p\to 0}r(p)=\infty,\quad \lim_{p\to 1}r(p)=0.$$
Thus,   the range of $r(p)$ is $[0,\infty)$. This proves that $u$ can be extended to $(0,\infty)$ satisfying the condition $u'(r)<1$ (item 4). From \eqref{p1}, let $p_R$ be the unique value such that 
$r(p_R)=R$. This yields 
\begin{equation}\label{c111}
c_1=R\frac{p_R}{(1-p_R^2)^2},
\end{equation}
where the value of $p_R$ needs to be determined.  Now \eqref{p3} becomes at $p=p_R$, 
$$u(p_R)=\frac{  R\, p_R}{(1-p_R^2)^2}\left(\frac14-\log(p_R)-p_R^2+\frac{3p_R^4}{4}\right).$$
As we require   $u(p_R)=M$, then  
\begin{equation}\label{c2}
\frac{p_R}{(1-p_R^2)^2}\left(\frac14-\log(p_R)-p_R^2+\frac{3p_R^4}{4}\right)=\frac{M}{R}.
\end{equation}
Define the function  
$$g\colon (0,1)\to\r,\quad g(x)=\frac{x}{(1-x^2)^2}\left(\frac14-\log(x)-x^2+\frac{3x^4}{4}\right).$$
Thus, the condition \eqref{c2} is fulfilled if we find a solution to 
\begin{equation}\label{c22}
g(x)=\frac{M}{R},\end{equation}
where we know $\frac{M}{R}\in (0,1)$.
The function $g(x)$ is strictly increasing with 
$$\lim_{x\to 0}g(x)=0,\quad \lim_{x\to 1}g(x)=1.$$
 This proves the existence of a unique solution  $p_R$ to \eqref{c22}.  Moreover, since $p=u'>0$, the function $u(r)$ is strictly increasing, proving the item (1).

 We prove (2). By the chain rule,  
$$\lim_{r\to 0}u'(r)=\lim_{p\to 1}\frac{u'(p)}{r'(p)}=\lim_{p\to 1}p= 1.$$

 
We study the regularity of $\gamma$. From \eqref{p1} and \eqref{p2}, we have   
\begin{equation}\label{p5}
\begin{split}
r'(p)&=-c_1(\frac{1}{p^2}+2-3p^2),\\
u'(p)&=-c_1p(\frac{1}{p^2}+2-3p^2).
\end{split}
\end{equation}
If   $\gamma'(p)=(0,0)$ for $p\not=0$,   the   second equation of \eqref{p5} gives  
\begin{equation}\label{p4}
\frac{1}{p^2}+2-3p^2=0.
\end{equation}
 This yields $p^2=1$. In such a case, $r(p)=0$, which  is outside the domain of $u$.   This proves that $\gamma$ is regular everywhere.

For the item (4), consider the expression for $u(p)$ given in \eqref{p3}. The function $r'(p)$ in \eqref{p5} is negative because $c_1>0$ and $1+2p^2-3p^4\geq 0$ if $p^2\leq 1$. Taking the variable $p$   for $u$, we have $u\colon [p_R,1)\to\r$. Since $r(p)$ in \eqref{p1} is defined for all $p\in (0,1]$, then $u$ can be extended up to $(0,1]$. Moreover, 
$$\lim_{r\to\infty}u(r)= \lim_{p\to 0}u(p)=\infty,$$
$$\lim_{r\to\infty}u'(r)= \lim_{p\to 0}\frac{u'(p)}{r'(p)}=0.$$ 

Finally, we prove that $u$ is concave, that is, $u''(r)\leq 0$ for all $r$. Using the chain rule, and the expressions for $r(p)$ and $u(p)$ in \eqref{p1} and \eqref{p2}, we have 
$$u''(r)=\frac{r'(p)u''(p)-r''(p)u'(p)}{r'(p)^3}=\frac{p^6}{c1}(-1-2p^2+3p^4)\leq 0$$
for all $p\in (0,1]$. This completes the proof.

\end{proof}

\end{document}
